%% ma: L12-B08-0105
%% lop: 12
%% chuong: 2
%% bai: 8
%% dang: 12.8.27
%% loai: TN4
%% muc: NB
%% dap_an: "A"
%% nguon: "CĐ2.3.pdf giáo viên gửi 27/10/2026 (M07, Lê Minh Khôi – Nguyễn Anh Khôi, Chương II Chủ đề 3 Biểu thức tọa độ của các phép toán vectơ), Câu 16"
%% kien_thuc: "Bài 8 (tọa độ trọng tâm tứ diện), Bài 6 (trọng tâm tứ diện)"
%% trang_thai: chinh-thuc
%% ghi_chu: "Biến thể thay số của L12-B08-0034"
\begin{ex}
Trong không gian $Oxyz$, cho tứ diện $ABCD$ có tọa độ các đỉnh $A(2;1;0)$, $B(3;7;4)$, $C(-2;1;-5)$ và $D(5;-9;3)$. Tọa độ trọng tâm $K$ của tứ diện $ABCD$ là
\choice{\True $K\left(2;0;\dfrac12\right)$}{$K\left(-2;0;-\dfrac12\right)$}{$K\left(2;0;-\dfrac12\right)$}{$K\left(-2;0;\dfrac12\right)$}
\loigiai{$K\left(\dfrac{2+3-2+5}4;\dfrac{1+7+1-9}4;\dfrac{0+4-5+3}4\right)=\left(2;0;\dfrac12\right)$.}
\end{ex}
